\documentclass{article}
\usepackage[utf8]{inputenc}
\usepackage[italian]{babel}
\usepackage{amsmath, amssymb}
\usepackage{tikz}

\begin{document}

\begin{flushright}
FISICA 1 - Fabrizio Cei \quad 05/03/2024
\end{flushright}

\vspace{0.5cm}
\noindent
\begin{minipage}{0.48\textwidth}
\begin{tikzpicture}[scale=0.8]
\draw[->] (-0.2,0) -- (5,0) node[right] {$\omega$};
\draw[->] (0,-0.2) -- (0,4) node[above] {$A$};
\node[left] at (0,3.5) {$A_M$};
\node[left] at (0,1) {$\frac{F_0}{\omega_0^2}$};
\draw[dashed] (0,3.5) -- (2.5,3.5);
\draw[dashed] (2.5,0) -- (2.5,3.5);
\node[below] at (2.5,0) {$\omega_0$};
\draw[thick, blue] (0,1) to[out=0,in=180] (2.5,3.5) to[out=0,in=170] (4.5,0.2);
\end{tikzpicture}
\end{minipage}
\hfill
\begin{minipage}{0.48\textwidth}
\begin{tikzpicture}[scale=0.8]
\draw[->] (-0.2,0) -- (5,0) node[right] {$\omega$};
\draw[->] (0,-0.2) -- (0,4) node[above] {$\varphi$};
\node[left] at (0,3) {$\pi$};
\node[left] at (0,1.5) {$\pi/2$};
\draw[dashed] (0,3) -- (4.5,3);
\draw[dashed] (0,1.5) -- (2.5,1.5);
\draw[dashed] (2.5,0) -- (2.5,1.5);
\node[below] at (2.5,0) {$\omega_0$};
\draw[thick, blue] (0,0) to[out=0,in=180] (2.5,1.5) to[out=0,in=180] (4.5,3);
\end{tikzpicture}
\end{minipage}

\vspace{0.5cm}
\noindent
\begin{minipage}[t]{0.48\textwidth}
\begin{align*}
A(\omega \ll \omega_0) &\cong \frac{F_0}{\omega_0^2} \\
A(\omega \gg \omega_0) &\cong \frac{F_0}{\omega^2} \to 0
\end{align*}
\end{minipage}
\hfill
\begin{minipage}[t]{0.48\textwidth}
\begin{align*}
\varphi(\omega \ll \omega_0) &\cong 0 \qquad \qquad 0 \le \varphi \le \pi \\
\varphi(\omega \gg \omega_0) &\to \pi
\end{align*}
\end{minipage}

\vspace{0.5cm}
\begin{center}
\begin{tikzpicture}
\node (A) at (0,0) {$A(\omega = \omega_0) = \frac{F_0}{\lambda \omega} = \frac{F_0 Q}{\omega^2}$};
\node (B) at (6,0) {$\varphi(\omega = \omega_0) = \pi/2$};
\node (R) at (3,-1) {\textbf{RISONANZA}};
\draw[->] (3,-0.2) -- (R);
\end{tikzpicture}
\end{center}

\vspace{0.5cm}
\begin{itemize}
\item \textbf{ENERGIA E RISONANZA}
\end{itemize}
Il periodo d'oscillazione della forzante è $T = 2\pi/\omega$.

Se si considera il SISTEMA OSCILLATORE + FORZANTE, è NON CONSERVATIVO perché l'energia trasmessa all'oscillatore per compensare quella persa è a spese della forzante.
Il SISTEMA OSCILLATORE invece è CONSERVATIVO poiché l'energia persa è recuperata grazie alla forzante.

\begin{equation*}
E = \frac{1}{2}m\dot{x}^2 + \frac{1}{2}kx^2 \qquad \frac{dE}{dt} = \dot{x}(m\ddot{x} + kx) = \dot{x}(-\beta\dot{x} + F_0\sin(\omega t))
\end{equation*}

\begin{equation*}
\left\langle \frac{dE}{dt} \right\rangle_T = \frac{1}{T} \int_t^{t+T} dt^* \left( \frac{dE}{dt^*} \right) = \frac{1}{T} [E(t+T) - E(t)] = 0 \quad \begin{subarray}{l} \text{moto periodico} \\ \text{perché } E(t+T)=E(t) \\ \text{Lo verifichiamo:} \end{subarray}
\end{equation*}

\begin{equation*}
= \frac{1}{T} \int_t^{t+T} dt^* \left[ \underbrace{-\beta(\omega A \cos(\omega t^* - \varphi))^2}_{\begin{subarray}{c} \text{DISSIPATA (MEDIA)} \\ \parallel \\ (-\beta \omega^2 A^2 / 2) \cdot T \end{subarray}} + \underbrace{F_0 \sin(\omega t^*) \omega A \cos(\omega t^* - \varphi)}_{\begin{subarray}{c} \text{ASSORBITA (MEDIA)} \\ \parallel \\ (+F_0 \omega A \sin\varphi / 2) \cdot T \end{subarray}} \right]
\end{equation*}

\begin{align*}
\begin{subarray}{c} \text{sostituisco } A \\ \text{e } \sin\varphi \text{ di} \\ \text{prima} \end{subarray} \quad &= \frac{\omega A}{2} (F_0 \sin\varphi - \beta \omega A) = \frac{\omega A m}{2} (f_0 \sin\varphi - \lambda \omega A) \\
&= \frac{F_0 \lambda \omega}{\sqrt{\dots}} - \frac{\omega \lambda F_0}{\sqrt{\dots}} = 0 \quad \checkmark \quad \text{l'oscillatore è conservativo}
\end{align*}

\vspace{0.5cm}
\begin{itemize}
\item Nel caso di $\varphi = \pi/2$ l'energia assorbita è $\frac{F_0 \omega A_M}{2}$ che è la massima possibile ($\sin\varphi, A$ sono MAX).
\end{itemize}
\begin{equation*}
\dot{x} = \omega A \cos(\omega t - \pi/2) = \omega A \sin(\omega t)
\end{equation*}
\begin{equation*}
\implies \frac{dE_{\text{ass}}}{dt} = F_0 \sin(\omega t) \cdot (\omega A \sin(\omega t)) = F_0 \omega A \sin^2(\omega t) \ge 0 \quad \forall t
\end{equation*}

È sempre LA SORGENTE che cede energia NELLA RISONANZA; nel resto dei casi ci sono dei momenti in cui anche la sorgente riceve energia dall'oscillatore.

\end{document}
